\clearpage
\SourceLeaf{148}
\begin{center}
\textsc{Gauss} an \textsc{Schumacher}. Göttingen, 6. Juli 1840.
\end{center}

\noindent . . . . . . . . . . .

Sie wissen, dass ich selbst auf das von mir seit 1794 gebrauchte Verfahren, dem später der Name Méthode des moindres quarrés beigelegt ist, niemals grossen Werth gelegt habe. Verstehen Sie mich recht; nicht in Beziehung auf den grossen Nutzen, den sie leistet, der ist klar genug, aber danach taxire ich die Dinge nicht. Sondern deshalb oder in so fern legte ich nicht viel Werth darauf, als vom ersten Anfang an der Gedanke mir so natürlich, so äusserst nahe liegend schien, dass ich nicht im Geringsten zweifelte, viele Personen, die mit Zahlenrechnung zu verkehren gehabt, müssten von selbst auf einen solchen Kunstgriff gekommen sein, und ihn gebraucht haben, ohne deswegen es der Mühe werth zu halten, viel Aufhebens von einer so natürlichen Sache zu machen. Namentlich fiel mir vor allen \textsc{Tobias Mayer} ein, und ich erinnere mich sehr bestimmt, dass ich \textls{oft}, wo ich mit andern von meiner Methode sprach (wie z.\,B.\ während meiner Studirzeit 1795--1798 wirklich vielfach geschehen ist), geäussert habe, ich wolle die allergrösste Wette eingehen, dass \textsc{Tobias Mayer} bei seinen Rechnungen dieselbe Methode schon gebraucht habe. Ich weiss nun jetzt aus jenen Papieren, dass ich jene Wette \textls{verloren} haben würde. In der That enthalten sie Eliminationen, z.\,B.\ von 3 unbekannten Grössen aus 4 oder 5 Gleichungen, aber so, wie es der ordinärste Rechner machen würde, ohne alle Spur irgend einer subtilern Kunst.

\noindent . . . . . . . . . . . .

